%\documentclass[english, 10pt,twocolumn]{article}
\documentclass[
reprint,
showkeys,
superscriptaddress,
aps,
pra,
longbibliography,
nofootinbib
]{revtex4-2}
\usepackage{placeins}	% for \FloatBarrier
%\usepackage{abstract}
%\usepackage{amsmath}
\usepackage{amssymb}
\usepackage{amstext}
\usepackage{amsthm}
\theoremstyle{definition}\newtheorem{definition}{Definition}
\usepackage[hyperref]{xcolor}
\usepackage{amsfonts}
\usepackage{mathrsfs}
\usepackage{mathtools}
\usepackage{dsfont}
\usepackage{eurosym}
%\usepackage[affil-it]{authblk}
%\usepackage[ansinew]{inputenc}
%\usepackage[latin1]{inputenc}

\usepackage[
colorlinks,
pdfpagelabels,
breaklinks,
pdfstartview=FitH,
bookmarksopen=true,
bookmarksnumbered=true,
bookmarksopenlevel=2,
plainpages=false,
hypertexnames=false,
pdftitle={Why Proof-of-Authority Consensus Is Unsuitable for Commercial-Bank-Operated DLT Infrastructures},
pdfauthor={Marc Winstel},
pdfkeywords={},
]{hyperref}
\usepackage{xurl}
\usepackage{bookmark}
\usepackage[capitalise,sort,compress]{cleveref}	% for clever references
% for oxford comma in cleveref
\newcommand{\creflastconjunction}{, and\nobreakspace}
\usepackage{url}
\usepackage{booktabs}
\usepackage{tabularx}


%\usepackage[backend=biber, style=chem-angew, sorting=none, sortcites=true, giveninits=true, maxbibnames=1, maxcitenames=4, giveninits=true, doi=false, eprint=true,url=true,natbib=true]{biblatex} 
%\addbibresource{Bib.bib}

\makeatletter
\renewcommand\paragraph{% new paragraph style
	\@startsection{paragraph}{4}{\z@}%
	{1.5ex plus 0.5ex minus 0.2ex}% space before
	{-1em}% space after; negative = run-in heading
	{\normalfont\normalsize\bfseries}% style
}
\makeatother
\setcounter{secnumdepth}{3} % no paragraph headings



\newcommand{\todo}[1]{\textcolor{red}{XX #1 XX}}
\newcommand{\missingCit}{\todo{Missing cite}}


\usepackage[acronym]{glossaries-extra}	% for acronyms
%\title{Why PoA consensus is not suitable for commerical-bank operated DLTs}%\title{Why Proof-of-Authority Is a Poor Consensus Choice for Commercial-Bank DLT Infrastructures}%\title{About the failure of PoA in DLTs for traditional finance}
\begin{document}
	
	% % % % % % % % % % % % % % % % % % % % % % % % % % % % % % % % % % % % % % % %
	% % % % % % % % % % % % % % % % % % % % % % % % % % % % % % % % % % % % % % % %
	
	% the title page
%	\title{Why Proof-of-Authority Consensus Is Unsuitable for Commercial-Bank-Operated DLT financial market infrastructures}
\title{Why Proof-of-Authority Consensus Is Ill-Suited for Commercial-Bank-Operated DLT-Based Financial Market Infrastructures}
	
	\author{Marc Winstel}
	\email{marc.winstel@bundesbank.de}
	\thanks{The views expressed in this paper are those of the author and do not necessarily coincide with those of the Deutsche Bundesbank or the Eurosystem.}
	\affiliation{
	    Deutsche Bundesbank, Wilhelm-Epstein-Straße 14, 60431 Frankfurt am Main, Germany.
	}
	
	\date{\today}
	
	\begin{abstract}
		Distributed ledger technology (DLT) has been widely proposed as a foundation for interbank settlement, cross-border payments, and tokenized asset markets in the future financial market infrastructures. 
		In this context, systems are typically permissioned and rely on Proof-of-Authority (PoA) or similar consensus mechanisms to achieve scalability and operational efficiency.
		
		 
		This paper identifies structural governance tensions in PoA-based, commercial-bank operated DLT systems by analyzing the role of PoA through the lens of institutional market structure and economic incentives for participating or operating commercial banks.	
		We argue that the institutional viability of authority-based consensus depends less on technical performance than on the identity, mandate, and incentives of the authority controlling transaction ordering, inclusion, and technical finality.
		PoA with a single validating entity introduces strategic dependence on a competing entity and creates incentives for market fragmentation, since large competitors have little reason to rely on infrastructure operated by a rival. 
		Multi-validator PoA distributes control, but exposes the infrastructure to governance stress, coordination failures and ambiguity over authoritative state, especially when operational governance of systems themselves become economically contested.
	%To contextualize these findings, we contrast PoA-based designs with Real-Time Gross Settlement (RTGS) systems. RTGS demonstrates that single-validator architectures can be viable when operated by a neutral authority issuing a risk-free settlement asset. The absence of comparable privately operated systems highlights that the limitations of PoA are not primarily technical, but institutional.
		%This suggests that PoA consensus does not provide a robust foundation for unified commercial-bank DLT infrastructures and that alternative architectures must explicitly account for governance, incentives, and legal finality.
		%However, PoA consensus may be viable in DLT under neutral governance, such as operation by technical service providers or central-banks, due to the lack of incentives to reorder transactions or exploit informational asymmetries. 

		The analysis suggests that PoA is not unsuitable for financial DLTs per se, but is ill-suited where validation authority is exercised by competitors in the market rather than by a neutral public or institutionally separated operator.
	\end{abstract}
	
	\keywords{
	    tokenization, proof-of-authority, consensus, blockchain, digital assets, smart contracts
	}
	
	\maketitle
	
	\tableofcontents
	
	% % % % % % % % % % % % % % % % % % % % % % % % % % % % % % % % % % % % % % % %


%%%%%%%%%%%%%%%%%%%%%%%%%%%%%%%%%%%%%%%%%%%%%%%%%%%%%%%%%%%%%%%%%%%
	
	

\section{Introduction}
\label{sec:introduction}

Distributed ledger technology (DLT), including the tokenization of digital and real-world assets, is widely considered to be a major component of the future financial market infrastructure, with potential applications in interbank settlement, cross-border payments, securities settlement, collateral management, and tokenized asset markets \cite{bisUnifiedLedger2025, oecdTokenisation2024,cpmiDlt2017, ecbDlt2016}. 
The major promise of DLT is the possibility of a shared programmable infrastructure capable of reducing reconciliation complexity between and within institutions, increasing automation through smart contracts, enabling atomic settlement of assets and payments, and potentially also operating continuously across jurisdictions, monetary zones and time zones.
DLT-based systems are often expected to reduce settlement latency, lower operational costs, improve transparency and auditability, and simplify interoperability between fragmented financial infrastructures \cite{cpmiDlt2017,bisUnifiedLedger2023}.
These expectations have led to extensive experimentation by central banks, commercial banks, and infrastructure providers in areas such as wholesale central bank digital currencies (wCBDC), tokenized asset markets, securities settlement and cross-border payments \cite{bisMbridge2022,bisUnifiedLedger2023}.

%\paragraph{DLT-based financial market projects.}
A number of central-bank-driven projects illustrate this development.
Initiatives such as \emph{Project Inthanon-LionRock}, its successor \emph{Project mBridge} and \emph{Project Cedar} investigate shared DLT infrastructures for cross-border payments and foreign-exchange settlement involving commercial and central banks \cite{inthanonLionrock2020, bisMbridge2022,nyicCedarAppendix}. 
More recently, \emph{Project Agora} has aimed to improve cross-border correspondent banking by integrating tokenised commercial bank deposits and tokenised wholesale central bank money on a shared programmable platform while preserving the existing two-tier monetary system \cite{bis_press_agora_2024, bis_bulletin87_2024}.
Moreover, a growing number of industry-led projects aim to establish shared DLT infrastructures. 
Examples include \emph{Fnality} \cite{fnality}, which explores tokenized settlement mechanisms backed by funds held at central banks, \emph{Partior} \cite{partior}, which focuses on cross-border interbank settlement and tokenized commercial-bank liabilities, and the \emph{Canton Network} \cite{canton_network_whitepaper_2024}, intended as a ``network of networks'', aiming to provide synchronization infrastructure for tokenized financial-market applications. 
Also, the Commercial Bank Money Token proof of concept considers the European Public Network, a public-permissioned EVM-based infrastructure using PoA consensus for financial-sector tokenisation use cases \cite{DKBDI2024CBMT}.
\todo{Look into deloitte workshop sources for more real world projects} 


Although these projects differ substantially in architecture and governance, they illustrate a common pattern: transaction validation, ordering, or synchronization is typically restricted to predefined authorized entities rather than opened to anonymous public participation.
This authority-based design choice is natural in regulated financial markets, but it creates a structural governance problem when authority over the ledger state is exercised by entities that are simultaneously competitors in the markets served by the infrastructure.

This discussion paper argues that Proof-of-Authority (PoA) and closely related authority-based consensus models are structurally problematic when used as the consensus layer of shared financial-market DLT infrastructures operated by commercial banks.
The problem is not primarily technical, i.e., does not necessarily regard common objectives for the design of consensus algorithms.
The central issue is who is allowed to decide the ordering, inclusion, and finalization of economically significant transactions.

The analysis distinguishes two canonical configurations.
In a single-validator setup, one commercial institution controls transaction ordering and validation, creating strategic dependence on a rival and incentives for market fragmentation.
In a multi-validator setup, control is distributed across a consortium of commercial institutions, but the system becomes exposed to governance stress, coordination failures, and ambiguity about authoritative state under disputed conditions.

\paragraph{Related literature.} % Here for related works on dlt as financial market infrastructure and PoA and on the necessity of wholesale CBDC
This paper relates to two strands of literature.
First, it contributes to the growing literature and policy discussion on DLT-based financial-market infrastructures, tokenisation, wholesale settlement, and programmable payment systems.
Central-bank, policy, and regulatory work has examined the potential use of DLT in payment, clearing and settlement, the settlement of tokenised assets in central-bank money, cross-border wholesale CBDC platforms, and unified-ledger or programmable-platform architectures \cite{cpmiDlt2017, auerMonnetShin2021,bisUnifiedLedger2023,bisTokenisation2024,oecdTokenisation2024,projectHelvetia,projectJura,bisMbridge2022,nyicCedarAppendix}.
Second, the paper builds on the technical work on distributed consensus, Byzantine fault tolerance and permissioned blockchains with particular focus on authority-based consensus mechanisms \cite{lamportByzantine1982,lamport1998paxos,ongaro2014raft,hyperledgerfabric,brown2016corda,szilagyi2017clique,moniz2020ibft,saltini2019ibft2}.
While this literature explains how distributed systems can agree on transaction validation, inclusion, ordering, and state replication, the present paper focuses on a different question, relating to relevant work on the future of financial market infrastructure \cite{mills2016dlt,cpmiDlt2017,oecdTokenisation2020,bisUnifiedLedger2023,bisUnifiedLedger2025,oecdTokenisation2024}: whether the institutional allocation of authority to provide consensus is compatible with shared financial-market infrastructure when the validators are themselves competitors in the market.
In this respect, the analysis also connects to work on transaction-ordering power, front-running and maximal extractable value, and to the financial-market-infrastructure literature on governance, legal basis, and settlement finality \cite{daian2020flashboys,eskandari2019transparent}.



\paragraph{Outline.}
The remainder of the paper is organized as follows.
\Cref{sec:authority-based-consensus} defines PoA and related authority-based consensus models in the context of financial-market DLTs.
Then, \cref{sec:failure-modes} analyzes single-validator and multi-validator configurations under commercial-bank operation.
After a discussion of neutral governance and Real time Gross Settlement (RTGS) systems as an institutional benchmark in \cref{sec:neutral-governance}, we conclude in \cref{sec:conclusions}.


\section{Authority-based consensus in financial-market DLTs}
\label{sec:authority-based-consensus}

Conceptually, a DLT consists of a set of distributed nodes, which are computers or systems operated by participating entities, maintaining a shared ledger state through a consensus mechanism or algorithm \cite{schneider1990implementing,nakamoto2008,cpmiDlt2017}. 
The consensus algorithm describes how the nodes agree on the validity and ordering of transactions and ensures that participants converge on a common representation of the system state despite distributed operation.
Consequently, the consensus mechanism effectively determines which transactions become part of the authoritative ledger state and under which conditions state transitions are finalized.
Examples of distributed consensus mechanisms include classical distributed systems algorithms such as Paxos \cite{lamport1998paxos} or Raft \cite{ongaro2014raft} as well as blockchain-specific approaches such as Proof-of-Work (PoW)\cite{nakamoto2008}, Proof-of-Stake (PoS) \cite{buterincasper2018PoS}, or Proof-of-Authority (PoA) (see, e.g., \cite{MANOLACHE2022580,szilagyi2017clique,woodPoA}). 


Public DLT systems often rely on open participation and decentralized consensus at the cost of scalability, settlement latency and computational efficiency.
In financial market applications of DLTs, participants are supposed to be legally identified entities, primarily commercial banks and central banks, regulated financial market infrastructures, and infrastructure providers.
Naturally, the permissioned access to institutional DLT systems leads to adoption of permissioned consensus mechanisms with predefined validators \cite{auerMonnetShin2021}, in stark contrast to public networks like Bitcoin \cite{nakamoto2008} or Ethereum \cite{buterin2014ethereum}, which originally popularized DLT.


%Consequently, the consensus problem differs fundamentally from that of the public blockchains: rather than coordinating anonymous and potentially adversarial actors\footnote{In typical blockchain and state machine replication literature, the security model of DLTs is often called Byzantine fault tolerance, see also the Byzantine Generals Problem described in \cite{lamportByzantine1982}.}, the system must primarily ensure consistent state replication and transaction ordering among a fixed set of identifiable institutions that may cooperate operationally while remaining economically and strategically distinct.
Consequently, the consensus problem differs fundamentally from that of the public blockchains: rather than coordinating anonymous and potentially adversarial actors, the system must primarily ensure consistent state replication and transaction ordering among a fixed set of identifiable institutions that may cooperate operationally while remaining economically and strategically distinct.
This distinction is central for financial-market DLTs.
The relevant trust problem is not eliminated by legal identification of validators, but transformed into a question of how ordering and validation authority is allocated among institutions with potentially divergent incentives.




\paragraph{Proof-of-Authority and centralized consensus.}

In this paper, PoA is used in a broad sense to denote permissioned consensus mechanisms in which transaction ordering, validation, or finality depends on a legally identified and administratively selected set of authorities.
Some variants of PoA also allow validators to vote about the introduction of new validators according to predetermined rules.
This includes explicit Ethereum-derived PoA protocols such as Clique, which is an Ethereum PoA protocol specified in EIP-225 \cite{szilagyi2017clique}, and Byzantine-fault-tolerant authority-based protocols such as IBFT and IBFT~2.0 \cite{moniz2020ibft,saltini2019ibft2}, but also functionally equivalent authority-based arrangements in commercial-bank DLT systems, such as notary services (see, e.g., \cite{cordaNotary}), ordering services (as in \cite{fabricOrdering}), validator committees, or synchronization or sequencing committees \cite{canton_network_whitepaper_2024}.
Enterprise Ethereum clients such as Hyperledger Besu and GoQuorum implement PoA-style protocols including Clique, IBFT~2.0, and QBFT for private permissioned networks \cite{besuPoa,goquorumConsensus}.

The purpose of this terminology is not to claim that all such systems implement the same protocol.
Rather, PoA is used as an analytical proxy for restricted consensus authority in systems where a selected set of legally identifiable institutions determines the authoritative ledger state.



The appeal of the PoA approach is clear. 
By eliminating open participation and competition from the consensus, PoA enables high throughput, low latency, and operational simplicity due to the reduction of the coordination problem, some of the most prominent drawbacks of early public DLT networks, while still providing advantages such as programmability and atomic settlement.
These technical advantages do not resolve the institutional allocation problem.
In financial-market infrastructures, the relevant question is not only whether a restricted validator set can process transactions efficiently, but whether the institutions controlling ordering and finality have incentives compatible with neutral market infrastructure.

Ultimately, the role of validators is economically significant, since they determine the ordering and finalization of payments, securities transfers, and other high-value financial transactions. 
Transaction ordering and inclusion is particularly relevant because it determines the underlying system state under which transactions are executed.
Thus, it directly affects economic outcome of transactions and may also result in certain operations to be delayed, reordered, or even rejected depending on the effects of the evolving ledger state.
The literature on front-running and maximal extractable value illustrates that control over ordering can itself become economically valuable \cite{daian2020flashboys,eskandari2019transparent}.



Thus, the governance structure underlying authority-based consensus in financial-market DLT systems remains a central institutional design problem.
While permissioned validation identifies who may participate in consensus, it does not by itself determine how validation authority should be allocated among participating institutions in the long term.
This question is not merely technical.
It reflects the economic structure of the underlying market and becomes more contested as the value of assets processed on the platform grows.



For the following analysis, two canonical configurations are distinguished.
In the first configuration, a single entity controls transaction validation and ordering. 
In the second configuration, validation is distributed across multiple institutions within a consortium, which either performs joint validator operation or employs selection rules, such as, e.g., rotating block production, for alternating validator selection.
These configurations do not exhaust all possible institutional designs, but they capture the main authority structures relevant for commercial-bank-operated PoA and related permissioned consensus systems.
It determines whether the system is governed by unilateral control or by collective agreement among competitors, and thus directly affects participation incentives, robustness of consensus, and the (legal) interpretation of ledger state.
The next section analyzes the structural failure modes associated with both configurations.
Rather than focusing on specific implementations, the analysis abstracts from existing projects to their underlying consensus structures and closely examines the effects of centralization in the economically relevant transaction ordering in commercial-bank DLT systems.





%%% outlook for rest of work
%This paper synthesizes existing concerns about permissioned DLT governance, validator incentives, legal finality, and market fragmentation into a focused critique of PoA and, more generally, of centralized consensus models for commercial-bank-operated shared financial infrastructures.
%Therefore, two setups of PoA are formalized and examined with respect to their implications in the context of commercial-bank DLT systems. 
%Rather than focusing on specific implementations, the analysis abstracts from existing projects to their underlying consensus structures and closely examines the effects of centralization in the economically relevant transaction ordering in commercial-bank DLT systems.


\section{Structural failure modes of shared market infrastructure under commercial-bank validation}
\label{sec:failure-modes}

The following analysis concerns DLTs with authority-based consensus layer operated by commercial banks and intended to function as shared financial-market infrastructures among multiple commercial banks.
The problem with PoA in commercial-bank-operated systems is not that the mechanism makes such systems technically unusable.
Rather, the problem is that its trust model does not match the institutional structure of the competition in financial markets.
A PoA ledger must answer a simple question: who is allowed to decide the canonical ordering of economically relevant transactions?
\Cref{tab:poa-configurations} summarizes the institutional incentive problems that arise from different allocations of validation authority in commercial-bank-operated DLT infrastructures.
\begin{table*}[ht]
	\centering
	\small
	\setlength{\tabcolsep}{8pt} % default is usually 6pt
	
	\begin{tabularx}{\textwidth}{p{0.16\textwidth} p{0.21\textwidth} p{0.21\textwidth} X}
		\toprule
		\textbf{Configuration} & \textbf{Authority structure} & \textbf{Immediate benefit} & \textbf{Institutional concern} \\
		\hline\hline \addlinespace
		Single commercial-bank validator 
		& One competing market participant controls transaction ordering and validation 
		& Operational simplicity, low coordination overhead, high performance 
		& Asymmetric dependence on a rival, potential censorship or reordering power, platform dependence for smaller participants, and fragmentation incentives for sufficiently large competitors \\
		\midrule
		Commercial-bank validator consortium 
		& Multiple competing institutions jointly determine the canonical ledger state 
		& Distribution of authority across participants, reduced unilateral control 
		& Coordination failures, governance disputes, fork or finality ambiguity, complex ex-post legal resolution \\
		\midrule
		Neutral governance through public validator or independent technical service provider
		& A non-commercial authority or independent service provider operates or governs the authoritative settlement layer 
		& Clear accountability, unified participation, legally supported finality 
		& Requires public mandate, legal basis, and institutional neutrality; not replicable by ordinary commercial-bank governance \\
		\hline\hline 
	\end{tabularx}
	\caption{Canonical authority-based consensus configurations in financial-market DLT infrastructures. The relevant distinction is not only the number of validators, but whether validation authority is exercised by a neutral public operator or by commercially competing market participants. The classification is specific to PoA and related authority-based consensus models and is not intended as a general taxonomy of all consensus mechanisms. The table does not imply that central-bank validation is necessarily appropriate for every conceivable shared financial-market DLT infrastructure.}
	\label{tab:poa-configurations}
\end{table*}



In commercial-bank infrastructures with multiple competitors, no neutral answer emerges naturally. 
If one of the competing firms controls validation, the system creates asymmetric dependence on that institution.
Small market participants may accept this dependence in exchange for access to network effects, liquidity, or operational infrastructure.
For sufficiently large competitors, however, such dependence creates incentives for non-participation or rival infrastructures.
If several competing entities jointly validate, the system must handle disagreement, faults, and governance disputes over the canonical chain. 
In both cases, the supposed simplicity of PoA reintroduces coordination and trust problems in institutions that consensus in DLTs shall originally resolve.
We analyze the two different cases in depth in the following.


\subsection{Single-validator Proof-of-Authority: centralization and market fragmentation}


In the single-validator case, PoA collapses into a centrally operated ledger. 
This may be acceptable when the validator is a central bank, public authority, or legally mandated infrastructure operator without being a competitor in the financial market.
It is much less plausible when the validator is a commercial bank.
A major bank is unlikely to rely on a ledger whose transaction ordering, inclusion, and operational availability are controlled by a direct competitor.

The problem is not solely reputational.
A single validator has privileged technical power over the ledger state.
It can delay transactions, censor transactions, reorder transactions, or observe commercially sensitive flows before they are finalized.
As discussed above, transaction ordering by itself can significantly affect the ledger state under which transactions are executed and can be exploited for the purpose of the validator, see, e.g., \cite{daian2020flashboys,eskandari2019transparent}.
Even if selfish reordering and further malicious behavior is legally prohibited, the technical possibility creates a strategic dependency on the validating institution.
The concern is particularly relevant in high-value financial-market infrastructures, where transaction ordering, inclusion, and timing can have direct economic consequences.


Legal accountability alone does not remove this problem.
Centrally operated blockchain systems can in principle achieve high performance and process large transaction volumes \cite{lovejoy2023hamilton,lovejoy2023parsec}.
In such settings, large amounts of economic value may depend on transaction ordering and inclusion decisions.
The volume and frequency of transactions would make it difficult for any jurisdiction to monitor, detect, and adjudicate manipulative behavior fully ex post.
Legal enforcement is therefore an important safeguard, but cannot substitute the neutral governance of the consensus layer itself.






Consequently, competing banks\footnote{At least, the competitors large enough to operate such a potentially huge and important system on their own.} would have strong incentives either not to join the platform or to create their own equivalent platforms. 
For sufficiently large competitors, the dependence on the validator is unlikely to be stable due to these incentives.
The likely result is market fragmentation: instead of converging on a limited number of shared infrastructures, the market splits into many bank-operated ledgers with limited interoperability and duplicated liquidity. 
In the extreme, each major competitor may operate its own DLT, turning the DLT infrastructure landscape into a set of rival settlement islands.
Interoperability would then have to be restored through bridges, APIs, contractual arrangements, or other reconciliation mechanisms.
Consequently, the operational complexity, liquidity fragmentation, and intermediation layers need to be reintroduced. 
These are precisely the frictions that shared DLT infrastructure promises to reduce.
In such a scenario, the DLT-specific value proposition becomes weak at the inter-institutional level.
The systems begin to resemble an updated form of internal replicated databases, with benefits concentrated mainly in interactions between dominant platform operators, the large market participants, and smaller participants.


\subsection{Multi-validator Proof-of-Authority: consortium governance and finality ambiguity}


The multi-validator case appears more attractive at first sight, since validation authority is distributed across several institutions rather than concentrated in a single entity.
In such a setup, the participating institutions may either jointly set up a validator consortium acting as the source of truth for the ledger, or validation may rotate among authorized institutions according to a predefined selection rule, such as round-robin alternation of block production.
Both variants reduce unilateral control.
However, they do not remove the institutional problem.
Rather, they transform the dependence on one commercial actor into collective governance among competing institutions, putting the agreement between the authorized validators on the terms of consensus operation, such as, e.g., validator selection, voting rights, upgrades, fees, access, emergency powers, under potentially enormous stress.



One might object that the many permissioned multi-validator systems are explicitly designed to prevent persistent forks, which are quite famous from the crypto community, for example by applying Byzantine-fault-tolerant protocols, quorum rules, or tightly governed validator consortia.
This objection is important, but it does not remove the institutional concern.
If a consortium of competing institutions is the ultimate source of truth, the system still requires a governance framework for validator admission, software upgrades, incident handling, rule changes, sanctions, emergency interventions, and dispute resolution.
The absence of routine forks is therefore not sufficient: the critical question is who controls the procedures that determine the authoritative state when the protocol, the participants, or the legal interpretation of ledger events come under stress.

It is noteworthy that validation through an inserted, validation-operating consortium may itself become a source of market fragmentation.
Different groups of institutions may form competing consortia, each claiming to provide a shared infrastructure for a segment of the market.
Such a structure may reduce unilateral dependence on one dominant bank, but it does not necessarily create a neutral and open infrastructure for the market as a whole.


Round-robin block production does not by itself imply forks, especially where proposer rotation is combined with quorum-based or Byzantine-fault-tolerant finality rules.
The institutional concern is different: when validators are also competing institutions, disputes over proposer selection, ordering rules, missed slots, sanctions, upgrades, or emergency interventions may become economically contested.
If the protocol or governance framework fails to resolve such disputes under stress, the result may be halted finality, emergency intervention, or, in extreme cases, competing claims about the authoritative ledger state.


A separately incorporated consortium or jointly governed operating entity may reduce some of these frictions by acting as the formal operator of the consensus layer on behalf of the participating institutions.
However, this does not eliminate the underlying governance problem.
The consortium remains shaped by the interests, voting rights, and strategic influence of its members.
Moreover, new entrants may face rules and technical choices that were established before they joined and that they cannot easily influence ex post.
Neutral and open market participation is, thus, prohibited.


The governance burden further increases with the number and heterogeneity of participants.
All changes to the DLT system's governance and operation including technical upgrades, validator rules, operational procedures, and emergency responses, must be coordinated among institutions with partly divergent commercial interests.
As the value settled on the platform grows, these governance decisions become economically more significant and therefore more contested.
Potentially, these conflicts may render a DLT system unsteerable and, in extreme cases, may lead to governance deadlock.
Legal frameworks for these systems must therefore specify not only normal operation, but also failure handling, validator disputes, emergency coordination, and the relationship between technical finality and legal finality.
This last issue is not unique to PoA: the relationship between technical finality and legal finality must be addressed for any programmable DLT system representing legally recognized assets.
Where such systems process high-value payments, securities, or tokenized claims, unresolved finality disputes may create spillovers beyond the platform itself and, in sufficiently large or interconnected infrastructures, may even become a source of systemic risk.


The relevant failure mode of multi-validator PoA is therefore not simple centralization, but institutional ambiguity.
Control is distributed, but the authority to resolve disagreement, define finality, and adapt the system remains embedded in a consortium of competing commercial actors.
For high-value financial-market infrastructures, this creates governance and legal-finality risks that cannot be solved by the consensus protocol alone.







\section{Neutral governance as the institutional condition for authority-based consensus}
\label{sec:neutral-governance}

The previous section analyzed the structural failure modes of PoA consensus when applied to shared market infrastructures operated by commercial banks or other competitors in the market.
This section argues that these failures are not inherent to authority-based consensus as such, but arise from the institutional position of the authority exercising control over transaction ordering, inclusion, and finality.
Thus, we examine how neutral governance makes authority-based consensus viable by demonstrating the mitigation of above failure modes through neutral authority consensus and discussing existing RTGS systems as institutional analogues for successful authority-based infrastructures under neutral governance.


%Thus, we examine how neutral governance makes authority-based consensus viable.
%First, we briefly demonstrate that the above failure modes are mitigated when the validators in PoA consensus are no competing market participants.
%Second, we discuss existing RTGS systems as institutional analogues for successful authority-based settlement infrastructures under neutral governance. 
%Finally, we argue that the absence of relevant privately operated RTGS systems supports the view that shared financial market infrastructures require governance by either public authorities, such as central banks, or independent technical service providers whose role is institutionally separated from competitive banking activity.

\subsection{Why neutrality changes the institutional problem of Proof-of-Authority}
The failure modes identified above do not arise from authority-based consensus as such, but from the institutional position of the authority exercising control over ordering, inclusion and finality. 
In a commercial-bank-operated PoA system, authority is assigned to actors that are themselves competitors in the market served by the infrastructure. 
This creates a structural conflict between the collective function of the infrastructure and the private incentives of its operators. 
Validator authority can therefore become strategically sensitive, as outlined above.


Neutral governance changes this institutional problem because the authority is separated from the competitive interests of the participants using the infrastructure. 
If the validating entity is a public authority, an independent operator, or another institutionally neutral body, the dependence on authority no longer takes the form of dependence on a market rival. 
This does not eliminate the need for accountability, operational resilience, or clear governance procedures. 
It does, however, remove the central source of strategic mistrust that makes commercial-bank-operated PoA unsuitable for shared financial market infrastructures. 
Under neutral governance, authority-based consensus can therefore function as a technical mechanism providing non-discriminatory ordering, inclusion, and finality within the infrastructure.
In such a setting, the potential benefits of DLT, including atomic settlement and programmability through smart contracts, can be realized in an ideal way.


In practice, such neutrality may be provided by central banks, whose public mandate and established role in settlement finality make them institutionally well suited to operate or govern authority-based financial market infrastructures, but whose direct involvement may limit flexibility, innovation speed, or private-sector experimentation.
Alternatively, neutrality may be provided by independent technical service providers, which can offer operational expertise and greater implementation flexibility, but only if their role is confined to technical operation and embedded in governance arrangements that prevent capture by, or preferential treatment of, individual market participants.
A technical outsourcing arrangement alone does not create neutrality if ultimate control remains with a subset of competing market participants.

 

\subsection{Real Time Gross Settlement systems: Authority-based settlement infrastructures with neutral governance}


RTGS systems illustrate that an authority-based settlement infrastructure can be compatible with financial-market infrastructure where the operator is neutral, has a public mandate, provides or settles in risk-free central bank money, and is supported by clear legal finality. 
RTGS systems, therefore, serve as an institutional benchmark rather than a technical analogy to DLT consensus.
They show that centralization of the authoritative settlement layer is not necessarily problematic in itself. 
The decisive condition is that the authority is not a commercial competitor in the markets served by the infrastructure.

The fact that existing RTGS systems are not DLT systems does not weaken this comparison, because the argument concerns the institutional allocation of settlement authority rather than the specific technical architecture used to record transactions.
Replacing the database layer of an RTGS-like system with DLT would not, by itself, compromise its institutional suitability, provided that authority over validation, ordering, and finality remained neutral, legally grounded, and separated from commercial competition.


The institutional suitability of RTGS systems can be summarized by the following features, each of which addresses one of the strategic vulnerabilities that arise when authority is assigned to competing market participants:
\begin{itemize}
	\setlength\itemsep{0.1em}
	\item \textbf{Neutral settlement authority:} The central bank is not a commercial competitor; it has no incentive to exploit transaction ordering, access decision or information asymmetries.
	
	\item \textbf{Legal finality:} Settlement takes place in central bank money and is supported by clear legal frameworks; this removes ambiguity about the authoritative settlement state.
	
	\item \textbf{System-wide infrastructure:} RTGS systems provide a common settlement layer for high-value payments and related financial-market activity. Their role as shared infrastructure reduces incentives for fragmentation into rival settlement systems.
	
	\item \textbf{Trusted governance:} Central bank mandate (monetary and financial stability, payment system efficiency) aligns with system-wide interests rather than private incentives.
	
	\item \textbf{Operational responsibility and accountability:} Authority over the infrastructure is matched by identifiable responsibility for system operation, risk management, incident handling, and crisis procedures.

	\item \textbf{Single authoritative state:} RTGS systems maintain one authoritative settlement record. This avoids the institutional equivalent of competing ledger states and removes the need for market participants to coordinate around rival versions of settlement history.\\
	
\end{itemize}


\subsection{The limited role of private operation in RTGS systems}


The role of private operation in RTGS systems is limited and institutionally qualified. Existing arrangements do not generally assign autonomous settlement authority to a commercial bank or a subset of competing commercial banks. Where private entities are involved, their role is typically embedded in a framework of central-bank settlement, public oversight, or delegated technical operation. The Swiss SIC system illustrates this distinction: it is operated by SIX Interbank Clearing Ltd on behalf of the Swiss National Bank (SNB), while the SNB remains the relevant public authority behind the system \cite{snb_sic_2024, snb_sic_web}. 
Similarly, the historical CHAPS arrangement in the United Kingdom involved private-sector governance of the high-value payment system, but this model was replaced in 2017 by direct operation of CHAPS by the Bank of England in order to strengthen end-to-end risk management and financial stability \cite{boe_chaps_direct_delivery_2017,boe_rtgs_chaps_intro_2025}.


This distinction is important for PoA-based DLT financial market infrastructures. 
Private technical operation does not create the same institutional problem as private settlement authority, provided that control over access, rule changes, ordering, finality, and crisis procedures remains neutral and accountable. 
By contrast, a PoA system operated or governed by commercial banks assigns the authoritative settlement function to actors with direct competitive interests in the market served by the infrastructure. 

The RTGS experience therefore supports the central claim of this section: authority-based consensus is not unsuitable for financial market infrastructures per se; it is unsuitable when the authority is itself a competing market participant or coalition of competing market participants.
Consequently, existing RTGS systems in established financial markets do not support autonomous settlement authority by competing commercial banks, but instead point to central-bank operation, such as in T2 \cite{ecb_t2_2026} or Fedwire Funds \cite{fedwire_funds_2024}, operated by the Eurosystem and Federal Reserve System, respectively, or institutionally qualified private operation under central-bank settlement, public oversight, or delegated technical mandates, as illustrated by SIC, which is operated by SIX on behalf of and under the supervision of the SNB \cite{snb_sic_2024}, and Lynx, which is operated by Payments Canada while relying on settlement accounts at the Bank of Canada \cite{canada_lynx_bylaw_2021}.
Major privately operated high-value payment systems such as CHIPS and EURO1 do not constitute autonomous private RTGS systems but instead rely on liquidity-saving netting or hybrid settlement mechanisms and remain anchored in central-bank money or central-bank-operated settlement infrastructures \cite{chips_tch,eba_euro1_key_features}.

\section{Conclusions}

\label{sec:conclusions}
The contribution of this paper is to shift the assessment of permissioned DLT consensus from technical performance alone to market structure, governance, and legal finality. 
In commercial-bank-operated financial-market infrastructures, the central design problem is not merely how consensus is reached, but who is institutionally authorized to determine the canonical ordering and final settlement state. 
Proof-of-Authority (PoA) consensus is therefore not only a technical design choice. 
It allocates control over transaction inclusion, ordering, and finality to an identifiable authority or group of authorities whose institutional position matters for the viability of the infrastructure.


\subsection{Summary}

This paper has argued that PoA consensus is structurally ill-suited for DLT-based financial-market infrastructures operated by market participants. 
The problem is not that PoA lacks technical performance: in permissioned environments, it can provide efficient transaction validation and operational simplicity. 
Its weakness lies in the institutional allocation of authority. 
When the validator role is assigned to commercial banks or consortia of commercial banks, consensus authority is allocated to institutions that are themselves market participants and potential competitors of those relying on the infrastructure.

This creates a structural dilemma. 
A single-validator PoA consensus design concentrates control over transaction inclusion, ordering, and finality in one commercial bank, thereby creating strategic dependence on a rival and encouraging market fragmentation. 
A multi-validator PoA design distributes this authority, but does not make it neutral; instead, it shifts the problem into the governance layer, where competing validators may face coordination failures, ambiguous accountability, and high-value disputes over protocol rules, validator conduct, emergency interventions, or the authoritative ledger state. 
The conclusion is therefore not that PoA cannot establish technical consensus, but that commercial-bank-operated PoA does not provide an institutionally robust basis for a fair, open, and legally reliable financial-market infrastructure.

\subsection{Discussion}

The analysis does not imply that PoA consensus is unsuitable for all financial DLT applications. 
Rather, it identifies the institutional conditions under which authority-based consensus can be viable. 
PoA may be appropriate where the validating authority is not itself a competing market participant, where governance is embedded in a public mandate, or where the ledger does not provide a common settlement infrastructure for rival institutions. 
This distinction explains why central-bank-operated or central-bank-led DLT arrangements can rely on authority-based consensus more plausibly than commercial-bank-operated systems. 
A central bank does not occupy the same strategic position as a commercial bank: it does not compete with settlement participants in the relevant market, is not exposed to the same incentives to exploit informational asymmetries, and can anchor finality in an established public-law and monetary framework.


The broader implication is that consensus mechanisms for financial-market infrastructures cannot be assessed independently of institutional design.
In high-value settlement systems, the consensus layer is not merely a technical device for ordering messages among known nodes. It allocates authority over inclusion, ordering, and finality, and therefore determines who can produce the authoritative state of the system.
DLT-based financial-market infrastructures therefore require a careful separation between technical validation and institutional authority. Efficient consensus is not sufficient; the infrastructure must also provide a governance structure capable of producing a neutral, legally reliable, and final settlement state.



One interesting implication of this analysis, that needs to be further explored, is that future DLT-based financial-market infrastructures may require hybrid architectures that separate public settlement authority from private innovation at higher layers.
The author suggests a design where a layer-one (L1) infrastructure operated or governed by a public authority could provide the neutral settlement, synchronization, and finality layer, while privately operated layer-two (L2) networks could offer application-specific functionality, market services, or operational experimentation.
This would allow interbank settlement of high-value obligations to remain anchored in a neutral institutional framework, while still enabling market participants to develop DLT-based services on top of that foundation on their respective L2s.
In a setting where central banks provide such an L1, the issuance and settlement of tokenized wholesale central bank money can also be explored.

\section{Acknowledgements}
The author acknowledges useful discussions with Jan Boensel on the implications of consensus mechanisms for financial market infrastructure.
Further, the author acknowledges useful collaboration with Jan Boensel, David-Alexandre Guiraud, and Andrea Tundis on related topics.
The author thanks the general directorate ``digital euro'' at Deutsche Bundesbank, specifically all members of department D\euro{} 1 for continuing discussions on wholesale CBDC, DLT and related topics.
The writing of this manuscript was supported using artificial intelligence tools, specifically \textit{ChatGPT} by \textit{OpenAI} \cite{chatgpt} and \textit{Claude} by \textit{Anthropic} \cite{claude}.
The author carefully reviewed, revised, and validated all generated content.


\FloatBarrier
\bibliography{Bib.bib}

\end{document}